\chapter{Facial Injury}

\section*{Definition and Overview}
Facial injuries are injuries to the face, including cuts, bruises, burns, dental injuries, and fractures of the facial bones. The face is exposed and vulnerable, and injuries to it are common from falls, sport, assaults, and road accidents. Because the face contains the eyes, nose, mouth, and airways, facial injuries require careful assessment. Most facial injuries are minor and heal well, but some affect vision, breathing, or appearance and need specialist treatment.

\section*{Epidemiology}
Facial injuries are common, particularly in young men, with sport, assaults, and road accidents the leading causes. In older people, falls are a frequent cause. Dental injuries are very common in children and in sport, and a significant proportion of facial injuries involve more than one area of the face. Most facial injuries heal without long-term problems, but they cause significant pain, time off work, and, in some cases, scarring.

\section*{Aetiology and Pathophysiology}
Facial injuries arise from direct force to the face.
\begin{itemize}
    \item \textbf{Causes:} falls, sporting injuries, assaults, motor vehicle accidents, and burns
    \item \textbf{Soft tissue injuries:} cuts, bruises, grazes, and burns of the skin and deeper tissues
    \item \textbf{Dental injuries:} chipped, loosened, or knocked-out teeth
    \item \textbf{Fractures:} breaks in the nose, cheekbones, eye sockets, or jaws
    \item \textbf{Eye injuries:} injury to the eye or its surrounding bones
    \item \textbf{Mechanism:} the force damages the skin, muscles, bones, or structures of the face
    \item \textbf{Effects:} pain, swelling, bleeding, and in serious cases problems with breathing, vision, and eating
\end{itemize}

\section*{Clinical Features}
The features depend on the type and severity of the injury.
\begin{itemize}
    \item Cuts, grazes, and bruising of the face
    \item Swelling and tenderness
    \item Bleeding from the nose, mouth, or wounds
    \item A knocked-out or chipped tooth
    \item Deformity, such as a flattened cheek or bent nose
    \item Difficulty opening the mouth or an abnormal bite
    \item Numbness of the cheek, lip, or chin
    \item Double vision, bruising around the eye, or a sunken eye
\end{itemize}

\section*{Differential Diagnosis}
Other conditions can be confused with facial injuries.
\begin{itemize}
    \item Swelling from allergy or infection, rather than trauma
    \item Bruising from blood disorders or medicines
    \item Bell's palsy and other causes of facial weakness
    \item Dental pain and infection
    \item Eye conditions causing redness or swelling
\end{itemize}

\section*{When to Seek Medical Care}
See a doctor for cuts that are deep or gaping, suspected fractures, dental injuries, or any facial injury with pain, swelling, or deformity. Urgent care is needed for serious injuries.

\textbf{Call 000 (or local emergency services) for:}
\begin{itemize}
    \item Difficulty breathing or signs of airway obstruction
    \item Loss of consciousness or confusion after a facial injury
    \item Sudden loss of vision or severe double vision
    \item Severe bleeding that does not stop
    \item A facial injury with a suspected neck or spinal injury
\end{itemize}

\section*{Diagnosis and Investigations}
Assessment identifies the extent of the injury.
\begin{itemize}
    \item \textbf{History:} how the injury happened and any loss of consciousness
    \item \textbf{Examination:} the wounds, face, eyes, mouth, and bite
    \item \textbf{Imaging:} X-rays or CT scans for suspected fractures
    \item \textbf{Eye assessment:} for injuries near the eye
    \item \textbf{Dental assessment:} for dental injuries
    \item \textbf{Head injury assessment:} when the injury involved the head
\end{itemize}

\section*{Management}
Management depends on the type of injury.
\begin{itemize}
    \item \textbf{Minor cuts and grazes:} cleaning, dressings, and sometimes glue or stitches
    \item \textbf{Wound care:} tetanus vaccination and antibiotics when needed
    \item \textbf{Dental injuries:} re-implantation of a knocked-out tooth, and dental review
    \item \textbf{Burns:} cooling and specialist burns care
    \item \textbf{Fractures:} observation, or surgery to realign and fix the bones
    \item \textbf{Swelling:} ice and elevation
    \item \textbf{Scar care:} good wound care and scar treatments after healing
\end{itemize}

\section*{Complications}
\begin{itemize}
    \item Infection of wounds
    \item Scarring and changes in appearance
    \item Nerve injury causing numbness or weakness
    \item Double vision or damage to the eye
    \item Misalignment of the teeth or bite after jaw injury
    \item The psychological impact of a facial injury
\end{itemize}

\section*{Prognosis and Outcomes}
Most facial injuries heal well, and with good wound care the scarring is minimal. Dental injuries often have good outcomes with prompt care, and fractures heal with proper treatment, usually returning the face to normal appearance and function. Serious injuries require specialist care, but outcomes are generally excellent.

\section*{Prevention and Self-Care}
\begin{itemize}
    \item Wear seatbelts and helmets
    \item Wear mouthguards and face protection in contact sports
    \item Use protective equipment at work
    \item Avoid alcohol and risk-taking behaviour
    \item Make the home safe to prevent falls
    \item Care for wounds promptly to minimise scarring
    \item Seek early dental care after dental injuries
\end{itemize}

\section*{Sources and Further Reading}
The following reputable sources provide further information for healthcare students and patients seeking reliable, current advice about facial injuries.
\begin{itemize}
    \item Australian Dental Association. \textit{Dental injury first aid.} https://www.ada.org.au/
    \item Healthdirect Australia. \textit{Facial injuries and cuts.} https://www.healthdirect.gov.au/
    \item St John Ambulance Australia. \textit{First aid for injuries.} https://stjohn.org.au/
    \item Australian and New Zealand Association of Oral and Maxillofacial Surgeons. \textit{Facial trauma information.} https://www.anzoms.org/
    \item NSW Health. \textit{Head and facial injury advice.} https://www.health.nsw.gov.au/
    \item MSD Manual. \textit{Facial trauma.} https://www.msdmanuals.com/
\end{itemize}
